\chapter{Introduction}

Mobile application and digital media usage is still on the rising\cite{perkins_internet_2018}.
Apps are one of many possible ways to interact with customers and consumer.
A web appearance on social media or a dedicated web page can also fulfill these needs.
The way, the user interacts with a web page is more shallow, because of lower expectations.
A user needs to log in most of the time.

A mobile application can enable user specific features without logging in, because the user can be identified with an identification token, that is generated, while the user installs the app.
It is possible to activate special features like geofencing, push notifications or local data storage, that is synchronized with the platform cloud services (e.g. Apple iCloud or Google Drive).

An mobile app integrated more deeply into the device, than a web app.
However, web browser already support a lot of the features by now, that originally are only available to native apps.
The current development indicates, that the differences will be further eliminated by the browser vendors.

\section{Problem definition}

Mobile applications have significant disadvantages, that push the development of web apps.
Changes have to be approved by the phone vendors, in order to be able to release update to the app.
This causes development delays and reduces the ability to react to events or customer requirements.

In addition, managing a lot of apps (> 20) requires a lot of work, since all of them have to be maintained.
New operating system features have to be implemented in all of the apps, which can cause inconsistencies.
All apps have to be tested in detail, to make sure that all apps implement the required change correctly.

Developers need deep system knowledge to develop the apps.
Every change has to go through one of each platform experts, so development is expensive.

\section{Motivation}

use of native mobile features, camera, background gps, beacons, push notification (with payload). caching, fast ui, good battery life.

no need to depend on third party frameworks, direct use of phone vendor's apis.

\section{Distinction}

This thesis will not describe the development experience nor the initial expense of the project. Neither will this thesis evaluate financial aspects and at which point, the project was or is worth the development costs.


\section{Structure}

\subsection{Research Question}

The thesis aims to answer the question ”Can a code interpreter be integrated in any domain?”. In addition, the development approach is compared to procedu- res, that are described in popular literature.

lorem ipsum